(i) Two systems only, fully observed and noise-free; Lorenz-96 has $D=10$, far from the high-dimensional regime in which LSTMs and large or parallel reservoirs were studied~\cite{vlachas2018data,vlachas2020backpropagation}. (ii) Five seeds in the main comparison and three in the sweeps. Paired tests with 3 seeds have little power, so several sweep statements rest on mean orderings only. (iii) All three forecasters are our own reimplementations, one small model per family, tuned with small grids of unequal size and with several optima on a grid edge (Sec.~\ref{sec:setup}); the test-seed sweeps show that neither the MLP row nor, possibly, the ESN row of Table~\ref{tab:main} is that family's best on Lorenz-96. A more thorough search (as in~\cite{racca2021robust}) could change the gaps, and on Lorenz-96 a larger reservoir or a longer-trained MLP changes them visibly. (iv) The MLP and GRU use the optimiser setting tuned at 5000 samples at every training length, so their curves in Fig.~\ref{fig:ntrain} mix a data effect with a fixed optimiser budget. (v) The GRU is the only recurrent baseline; no LSTM, no next-generation reservoir~\cite{gauthier2021next}, no hybrid or physics-informed model. (vi) VPT uses one error threshold (0.4). The climate metric is a per-component marginal Wasserstein distance with one pass threshold; it does not check Lyapunov exponents or correlation structure, so ``attractor statistics'' means marginal statistics only, and the metric may be too coarse to separate the three systems. (vii) No experiment isolates a mechanism (why squared features or input noise stabilise long rollouts, why a small $\rho$ helps on Lorenz-96). (viii) This is the second version of the experiments: an audit of the first found an under-trained MLP, a non-deterministic ESN eigensolver and a noisier climate metric; all groups were rerun, and the earlier rows remain in the registry, marked superseded.
